\documentclass[a4paper, serif, 10pt]{moderncv}

%% moderncv
% style and color
\moderncvstyle{banking}
\moderncvcolor{black}

% renew moderncv command to adapt for CJK colon
\renewcommand*{\cvitem}[3][.25em]{%
  \ifstrempty{#2}{}{\hintstyle{#2}：}{#3}%
  \par\addvspace{#1}}

\renewcommand*{\cvitemwithcomment}[4][.25em]{%
  \savebox{\cvitemwithcommentmainbox}{\ifstrempty{#2}{}{\hintstyle{#2}：}#3}%
  \setlength{\cvitemwithcommentmainlength}{\widthof{\usebox{\cvitemwithcommentmainbox}}}%
  \setlength{\cvitemwithcommentcommentlength}{\maincolumnwidth-\separatorcolumnwidth-\cvitemwithcommentmainlength}%
  \begin{minipage}[t]{\cvitemwithcommentmainlength}\usebox{\cvitemwithcommentmainbox}\end{minipage}%
  \hfill% fill of \separatorcolumnwidth
  \begin{minipage}[t]{\cvitemwithcommentcommentlength}\raggedleft\small\itshape#4\end{minipage}%
  \par\addvspace{#1}}

%% page layout/margins
\usepackage[top=2.5cm, bottom=2.5cm, left=1.5cm, right=1.5cm]{geometry}

\nopagenumbers{}

%% Babel config for English language
\usepackage[english]{babel}

%% fontspec
\usepackage{fontspec}

\IfFontExistsTF{Linux Libertine}{
  \setmainfont[Ligatures={TeX, Common}, Numbers=Lining, ItalicFont=Linux Libertine]{Linux Libertine}
}{}
\IfFontExistsTF{Linux Libertine O}{
  \setmainfont[Ligatures={TeX, Common}, Numbers=Lining, ItalicFont=Linux Libertine O]{Linux Libertine O}
}{}

%% CTeX
% CJK support, used to show CJK characters in the resume
%
% - fontset=none: disable builtin fontset but instead set the CJK font manually
% - heading=false: disable ctex heading
% - punct=kaiming: use kaiming punctuations styles for CJK
% - scheme=plain: use plain scheme, do not override \normalsize font size
% - space=auto: space settings for CJK characters
%
% ref:
% - http://ctan.mirrorcatalogs.com/language/chinese/ctex/ctex.pdf
\usepackage[UTF8, heading=false, punct=kaiming, scheme=plain, space=auto]{ctex}

\IfFontExistsTF{Noto Serif CJK SC}{
  \setCJKmainfont{Noto Serif CJK SC}
}{}
\IfFontExistsTF{Noto Sans CJK SC}{
  \setCJKsansfont{Noto Sans CJK SC}
}{}

%% line spacing
\usepackage{setspace}
\setstretch{1.125}

%% URL styling - use normal text instead of monospace
\urlstyle{same}

% Auto-underline all links
% Defer until \begin{document} so hyperref is already loaded
\AtBeginDocument{%
  \let\oldhref\href
  \renewcommand{\href}[2]{\oldhref{#1}{\underline{#2}}}%
}

%% Patch \httplink and \httpslink to handle full URLs
% The original moderncv macros always prepend http:// or https:// to URLs.
% This causes issues when the URL already has a protocol (e.g., https://example.com)
% resulting in malformed URLs like https://https://example.com.
% This patch checks if the URL already contains "://" and uses it directly if so.
% Note: We use \str_set:Nx (x-expansion) to expand macros like \@homepage first.
\ExplSyntaxOn
\str_new:N \g_yamlresume_url_str

\RenewDocumentCommand{\httpslink}{O{}m}{
  \str_set:Nx \g_yamlresume_url_str {#2}
  \str_if_in:NnTF \g_yamlresume_url_str {://}
    { \url{#2} }
    { \url{https://#2} }
}

\RenewDocumentCommand{\httplink}{O{}m}{
  \str_set:Nx \g_yamlresume_url_str {#2}
  \str_if_in:NnTF \g_yamlresume_url_str {://}
    { \url{#2} }
    { \url{http://#2} }
}
\ExplSyntaxOff

\name{林靜文}{}
\title{資深客戶經理 | B2B SaaS 業務開發專家}
\phone[mobile]{+886 912 345 678}
\email{chingwen.lin@example.com}
\homepage{https://www.chingwenlin.tw}

\address{中國臺灣，台北市，台北市，台北市松山區敦化北路 123 號 8 樓，105}{}{}

\extrainfo{{\small \faLinkedin}\ \href{https://www.linkedin.com/in/chingwenlin}{@chingwenlin} {} {} {} • {} {} {} 
{\small \faLine}\ \href{https://line.me/ti/p/\textasciitilde{}chingwenlin}{@chingwenlin}}

\begin{document}

\maketitle

\section{簡介}

\cvline{}{具備 7 年 B2B SaaS 與數位解決方案銷售經驗的資深客戶經理，專注於企業客戶開發與關係維護。
%
\begin{itemize}
\item 累積管理超過 50 家企業客戶，年度業績配額達成率平均 125\%
\item 擅長顧問式銷售、合約談判與跨部門專案協調，能將客戶需求轉化為可落地的解決方案
\item 熟悉 CRM 系統與數據分析，透過客戶旅程優化提升續約率與客戶終身價值
\item 具備優秀的簡報與溝通能力，能與 C-level 決策者建立長期信任關係
\end{itemize}}
\section{教育背景}

\cventry{2015年9月–2019年6月}
        {學士，企業管理學系，成績：3.9/4.0}
        {國立臺灣大學}
        {\url{https://www.ntu.edu.tw}}
        {}
        {\begin{itemize}
\item 主修企業管理，輔修經濟學，奠定商業策略與數據分析基礎
\item 擔任系學會公關長，籌辦超過 10 場企業參訪與講座
\item 畢業專題「訂閱制商業模式之消費者留存分析」獲選為年度優良報告
\end{itemize}
\textbf{課程}：行銷管理、消費者行為、策略管理、財務管理、商業談判、資料庫管理}
\section{工作經歷}

\cventry{2022年1月至今}
        {資深客戶經理}
        {智雲科技股份有限公司}
        {\url{https://www.cloudthink.com.tw}}
        {}
        {\begin{itemize}
\item 負責台灣與大中華區企業客戶之 SaaS 解決方案銷售，管理年度 3,000 萬元業績配額
\item 連續兩年達成 120\% 以上業績目標，開發 20 家以上新客戶，包含 3 家財星 500 大企業
\item 主導客戶成功計畫，提升續約率至 92\%，客戶終身價值成長 35\%
\item 與產品、行銷及技術團隊協作，將客戶回饋轉化為產品優化建議
\end{itemize}
\textbf{關鍵字}：B2B SaaS、企業銷售、客戶關係管理、合約談判、業績成長}

\cventry{2019年7月–2021年12月}
        {客戶經理}
        {網訊科技股份有限公司}
        {\url{https://www.netinfo.com.tw}}
        {}
        {\begin{itemize}
\item 負責數位行銷與雲端服務解決方案之客戶開發與維護，涵蓋零售、金融與製造業
\item 達成季度業績目標 110\%，兩年內將負責區域營收提升 40\%
\item 建立標準化客戶拜訪與需求分析流程，縮短銷售週期 20\%
\item 協調內部資源處理客戶問題，客戶滿意度維持 4.8/5.0
\end{itemize}
\textbf{關鍵字}：數位行銷、雲端服務、業務開發、客戶滿意度}
\section{語言}

\cvline{中文}{母語或雙語水平 \hfill \textbf{關鍵字}：母語}
\cvline{英語}{完全專業水平 \hfill \textbf{關鍵字}：TOEIC 950、商務談判}
\cvline{日語}{有限工作水平 \hfill \textbf{關鍵字}：JLPT N2、日常會話}
\section{專業技能}

\cvline{客戶關係管理}{專家 \hfill \textbf{關鍵字}：CRM、Salesforce、HubSpot、客戶旅程優化}
\cvline{業務開發}{高級 \hfill \textbf{關鍵字}：陌生開發、顧問式銷售、通路開發、業績預測}
\cvline{簡報與談判}{高級 \hfill \textbf{關鍵字}：高階簡報、合約談判、跨部門溝通、衝突解決}
\cvline{數據分析}{中級 \hfill \textbf{關鍵字}：Excel、Tableau、SQL、Google Analytics}
\section{榮譽}

\cventry{2024年1月}
        {智雲科技股份有限公司}
        {年度最佳業務獎}
        {}
        {}
        {表揚年度業績達成率最高且客戶滿意度優異之業務人員。}

\cventry{2022年3月}
        {網訊科技股份有限公司}
        {卓越客戶服務獎}
        {}
        {}
        {表彰在客戶關係維護與服務創新方面有傑出表現的員工。}
\section{證書}

\cventry{2023年5月}
        {Salesforce}
        {Salesforce Certified Sales Cloud Consultant}
        {\url{https://trailhead.salesforce.com/credentials/salescloudconsultant}}
        {}
        {}

\cventry{2021年8月}
        {Google}
        {Google Analytics Individual Qualification}
        {\url{https://skillshop.exceedlms.com/student/catalog/list?category\_ids=53-google-analytics}}
        {}
        {}
\section{出版刊物}

\cventry{2023年11月}
        {顧問式銷售在 SaaS 產業的應用與實踐}
        {商業周刊}
        {\url{https://www.businessweekly.com.tw}}
        {}
        {\begin{itemize}
\item 探討顧問式銷售如何協助 B2B SaaS 企業提升客戶留存與擴張
\item 分享實務案例與客戶旅程優化框架
\end{itemize}}
\section{項目}

\cventry{2023年3月–2023年9月}
        {運用 CRM 數據與機器學習建立客戶流失預警模型}
        {企業客戶健康度預警系統}
        {\url{https://github.com/chingwenlin/customer-health}}
        {}
        {\begin{itemize}
\item 整合 Salesforce 與內部數據，建立客戶健康度指標
\item 透過 Python 與 SQL 分析歷史續約資料，識別高風險客戶
\item 系統上線後協助業務團隊提前介入，降低流失率 15\%
\end{itemize}
\textbf{關鍵字}：CRM、Python、數據分析、客戶成功}
\section{興趣愛好}

\cvline{運動}{馬拉松、羽球、登山}
\cvline{閱讀}{商業策略、心理學、人物傳記}
\section{志願服務}

\cventry{2021年9月至今}
        {業界導師}
        {台灣數位行銷協會}
        {\url{https://www.tdma.org.tw}}
        {}
        {\begin{itemize}
\item 擔任業界導師，協助年輕學子了解 B2B 業務職涯發展
\item 每季舉辦工作坊，分享客戶開發與簡報技巧
\item 參與協會產學合作計畫，協助媒合實習機會
\end{itemize}}
    
\end{document}